\chapter{Background and Related Work}
\label{ch:related}

This chapter (1) introduces the web-agent task formulation and the benchmark
used throughout the thesis, then reviews related work on (2) LLM-based web
agents, (3) skill and workflow discovery, (4) automatic training-data
generation for agents, and (5) small language models as agents, before
(6) positioning this thesis against these lines of work.

% Writing rule: group by topic, oldest to newest inside each topic, 2-4
% sentences per work, and END EVERY ITEM WITH A CRITIQUE ("However, ...").
% Every citation must pass the verification workflow (memory file 04).

\section{Background}
\label{sec:rel-background}
% TODO: web agent formulation (observation = accessibility tree, action
% space); WebArena and BrowserGym; ReAct-style agent loops; SFT and LoRA.
\textit{[To be written — see outline 07, section 2.1.]}

\section{LLM-Based Web Agents}
\label{sec:rel-llm-agents}
\textit{[To be written — see outline 07, section 2.2.]}

\section{Skill and Workflow Discovery}
\label{sec:rel-skills}
% TODO: options framework; Agent Workflow Memory; skill induction lines.
\textit{[To be written — see outline 07, section 2.3.]}

\section{Automatic Data Generation for Agent Training}
\label{sec:rel-datagen}
% TODO: survey and VERIFY candidates (NNetNav, OS-Genesis, Synatra, ...).
\textit{[To be written — see outline 07, section 2.4.]}

\section{Small Language Models as Agents}
\label{sec:rel-slm}
\textit{[To be written — see outline 07, section 2.5.]}

\section{Positioning of This Thesis}
\label{sec:rel-positioning}
% TODO: comparison table over axes: graph-structured exploration? verified
% tasks? long-horizon composition? trains an SLM?
\textit{[To be written — see outline 07, section 2.6.]}
